Haus- und Studienabschlussarbeiten werden im Format DIN~A4 angefertigt
(\qtyproduct{210 x 297}{\milli\metre}). Das Papierformat ist durch die
Klassenoption \texttt{a4paper} in der Vorlage bereits voreingestellt (siehe
Listing~\ref{lst:tpl:thesis}).

Ausnahmen gelten für Dissertationen und Habilitationen, die zugleich im
Buchverlag veröffentlicht werden (müssen). Wenn Sie also keine Dissertation
oder Habilitation verfassen, überspringen Sie den Rest dieses Abschnitts.

Buchveröffentlichungen erscheinen in einem kleineren Format. Die genauen
Abmessungen variieren je nach Verleger. Der Universitätsverlag benutzt für
Schriftenreihen i.\,d.\,R. das Format DIN~A5 (\qtyproduct{210 x 297}{\milli\metre}).
Autoren von solcher Arbeiten stehen nun vor einem Problem. Gutachter verlangen
die Arbeit meist im Format DIN~A4, die Veröffentlichung findet aber in einem
anderen Format statt. Es gibt zwei Lösungsmöglichkeiten, die beide ihre Vor-
und Nachteile aufweisen.

\begin{enumerate}
\item Sie legen die Arbeit für die Gutachter im Format DIN~A4 an, und im
      Rahmen der Veröffentlichung wird das Manuskript auf das Buchformat
      umgearbeitet. Dies bedeutet zusätzlichen Aufwand, geht aber gewissen
      Problemen aus dem Weg, die sich beim Druck ergeben können.

\item Sie reichen das Werk im Format DIN~A4 beim Drucker ein, der es für Sie
      auf das korrekte Format herunterskaliert. Sie müssen Ihre Arbeit dann
      nicht noch einmal überarbeiten. Es ist jedoch von großer Wichtigkeit,
      diesen Weg bereits bei der Erstellung der A4-Vorlage (mit anderen Worten
      von Anfang an) zu berücksichtigen.
\end{enumerate}

Entscheiden Sie sich für den ersten Weg, empfehlen wir Ihnen den Vorspann für
die A5-Version wie in Listing~\ref{lst:tpl:papersize_a5} dargestellt
anzupassen. Verständigen Sie sich vor dem Druck mit Ihrem Verlag bitte
unbedingt über den Wert der Bindekorrektur (siehe Abschnitt~\ref{sec:tpl:bcor}).

\lstinputlisting[language={[LaTeX]{TeX}},style=block,float,caption={Umstellung auf DIN~A5},label={lst:tpl:papersize_a5}]{code/tpl_papersize_a5.tex}

Wollen Sie den weniger aufwendigen zweiten Weg beschreiten, wählen Sie für die
A4~Vorlage mindestens die Basisschriftgröße~\qty{12}{\pt}. Bei der Skalierung
von A4 auf A5 wird das Dokument auf ca.~\qty{71}{\percent} seiner
Ursprungsgröße verkleinert. Ein Schriftzeichen der Größe~\qty{12}{\pt} ist dann
nur noch knapp \qty{9}{\pt} groß. \LaTeX\ verwendet bei vielen Schriftfamilien
jedoch für jede Schriftgröße eine eigene Schrift. Ein von \qty{12}{\pt} auf
\qty{9}{\pt} verkleinertes Zeichen hat damit ein anderes Erscheinungsbild als
ein Zeichen, welches für \qty{9}{\pt} entworfen wurden. Beim Herunterskalieren
sind die Linien oft dünner als in der jeweiligen Entwurfsgröße. Die Schriftart
"`Latin Modern"' ist hierfür besonders anfällig (siehe
Abschnitt~\ref{sec:tpl:font}).

Ist die Basisgröße in der DIN~A4-Vorlage zu klein gewählt, werden die
Schriftzeichen in der A5-Variante zu mager, und dünnere Linien können u.\,U.
sogar ganz verschwinden. Der Verlag wird Sie beim Probedruck sicher darauf
hinweisen, dass Sie Schriften außerhalb ihrer Entwurfsgröße verwenden und dass
er dadurch keinen optimalen Druck garantieren kann. Diese Warnung müssen Sie
dann hinnehmen.
